\section{Conclusion}
\label{sec:conclusion}

This paper put Vibe Theory under simulation and assembled the results into one picture. From the framework's two commitments, that reality is one substance and that the substrate is finite, we built a reproducible testbed and asked nine load-bearing questions, then measured the answers.

The headline is twofold. Monism is computationally coherent: geometry, time, matter, force, and the onset of gravity and quantum behaviour all come out of one mesh of tones and notes, with nothing added, as the dictionary of Section~\ref{sec:picture} records. And the dynamics does more than describe it: the substrate's own action, the discrete gravitational action, makes a smooth low-dimensional spacetime a stable phase that decays without the action and persists with it, coexisting with the layered phase in a first-order signature. So the smooth world is supported by the law, not merely assumed. Whether that phase strictly dominates is the open refinement. This is movement from a coherent description toward a generative theory for the physical ladder, a direction rather than a completed crossing.

The picture is also tight. The hyperbolic, expanding, no-hidden-state commitments are not independent style choices. They reinforce one another: the geometry that makes the substrate navigable is the geometry whose smooth phase the gravitational action stabilizes, the topology that gives the fermion its spin is what the gauge field couples to, and the connectedness that makes the mesh one structure is what allows the quantum correlation. The commitments stand or fall as a package, and so far they stand.

We have been equally precise about the limits. Which phase strictly dominates in the continuum is strong evidence, not proof. The quantum link is made quantitative, not closed: a natural mesh's correlation decays with separation while quantum violation does not, and that tension is the framework's most contested rung. The Weyl-projected chiral gauge theory, the full field equations of gravity, the specific Standard Model content, and the foundational claim about experience are open in physics or lie at the framework's boundary, and we marked each as such rather than overclaim.

What Vibe Theory offers, then, is a candidate ontology that has been simulated rather than only argued. Its physical ladder is built and largely validated, its hardest dynamical question is advanced to a stable phase and a first-order signature at the scale we can reach with strict dominance still open, its quantum question is sharpened to a precise open problem, and its boundary is named honestly. The substrate is one kind of thing, related to itself in different patterns, and the testbed shows that a great deal of the world really does come out of it. The next steps are clear: the free-energy crossing, the natural alignment of the quantum correlation, and the climb past the chiral wall. They are stated precisely enough to be attempted, simulated, and criticized, which is the most a model at this stage can ask.
